% =====================================================================
% Experiment 1: CMO-RL vs. CHEHAB RL + rotation key selection
% ---------------------------------------------------------------------
% Required packages (add to the preamble if not already loaded):
%   \usepackage{booktabs,xcolor,pgfplots,amssymb,pifont,cleveref}
%   \usepgfplotslibrary{groupplots}
%   \pgfplotsset{compat=1.17}
%
% ALL NUMBERS IN RED ARE ILLUSTRATIVE PLACEHOLDERS, NOT RESULTS.
% Blue [TODO] notes are open questions for the authors.
% Remove the two macros below (or redefine \ex{#1} as #1) once real
% results replace the placeholders.
% =====================================================================

% --- Macros (move to the preamble; keep names consistent with the paper)
\providecommand{\cmorl}{\textsc{CMO-RL}}
\providecommand{\chehabrl}{\textsc{CHEHAB~RL}}
\providecommand{\rks}{\textsc{RKS}}                 % rotation key selection
\providecommand{\baseline}{\chehabrl{}+\rks{}}
\providecommand{\wexec}{w_{\mathit{exec}}}
\providecommand{\wkeys}{w_{\mathit{keys}}}
\providecommand{\budget}{\mathcal{B}}
\providecommand{\ex}[1]{\textcolor{red}{#1}}        % placeholder value
\providecommand{\todo}[1]{\textcolor{blue}{[TODO: #1]}}

\subsection{Experiment 1: \cmorl{} versus \chehabrl{} with Rotation Key Selection}
\label{sec:exp1}

This experiment asks whether a single \cmorl{} policy can satisfy a noise
budget that the original compiler violates, without giving up the latency
that the original compiler achieves. We compare \cmorl{} with the original
\chehabrl{} optimizer~\cite{chehabrl} followed by its rotation key selection
pass. We fix one noise budget and one preference vector, compile every
benchmark with both systems, and measure three quantities for each compiled
circuit: the noise it consumes, its execution time, and the size of the
rotation keys it requires.

\paragraph{Baseline.}
\chehabrl{} rewrites a program by applying term rewriting rules selected by a
PPO agent. Its reward is derived from a single cost function that combines
operation cost, circuit depth, and multiplicative depth. This cost function
targets execution time only. \chehabrl{} is therefore unaware of two
properties that matter in deployment: the number of distinct rotation keys
that the circuit requires and the noise that the circuit accumulates. To
control the first property, \chehabrl{} applies a post-optimization pass that
we call rotation key selection (\rks{}). \rks{} counts how often each
rotation offset is used and rewrites infrequent rotations as chains of
power-of-two rotations using their non-adjacent form (NAF). For example, a
rotation by~12 becomes a rotation by~8 followed by a rotation by~4, which
removes the dedicated key for offset~12 when the keys for offsets~8 and~4 are
already present. \rks{} stops once the number of keys falls below a fixed
threshold. \rks{} reduces key memory, but it runs after the agent has fixed
the structure of the circuit, and it does not consider noise. We refer to the
complete baseline as \baseline{}.
\todo{State the \rks{} key threshold used here and how it was chosen. For a
fair comparison we suggest either the default threshold of the original
implementation or a threshold equal to the key count produced by \cmorl{}.}

\paragraph{\cmorl{}.}
\cmorl{} is a single policy conditioned on two inputs: a preference vector
$\vec{w} = (\wexec, \wkeys)$, which sets the relative weight of execution
time and rotation key memory, and a noise budget $\budget$, which bounds the
noise that the compiled circuit may consume (\cref{sec:method}). The policy
chooses rewrites for both objectives jointly. It can therefore select a
vectorization strategy that uses fewer rotation offsets, instead of
decomposing rotations after the strategy has been fixed. We evaluate the raw
output of the policy, with the safety rollback described in
\cref{sec:rollback} disabled. As a result, every constraint satisfaction we
report comes from the policy itself and not from the rollback mechanism.
\todo{Confirm that rollback is disabled for Exp.~1; if it is enabled, say so
and report how often it was triggered.}

\paragraph{Setup.}
We set the noise budget to $\budget = \ex{230}$~bits of consumed BFV noise.
We chose this value because it lies in the range where the circuits produced
by \baseline{} exceed the budget for a majority of the benchmarks, while
satisfying it is still possible: it is the tightest of the budgets that
\cmorl{} was trained on.
\todo{Confirm the budget value and its meaning (consumed noise bits as
estimated by the noise model vs.\ the remaining noise budget measured by SEAL
for fixed encryption parameters $(N, \log q)$). State the encryption
parameters.}
We set the preference vector to $\vec{w} = (\wexec, \wkeys) = (0.9, 0.1)$.
This preference keeps execution time as the primary objective while assigning
a non-zero weight to rotation key memory; it is the preference closest to the
single objective of \chehabrl{} that still asks \cmorl{} to co-optimize
memory. \baseline{} has no preference input and is run unchanged.
We evaluate both systems on the \ex{14} benchmarks listed in
\cref{tab:exp1-details}, none of which appears in the training set.
\todo{Final benchmark list; the list below is a placeholder subset of the
thesis benchmarks.}
For every compiled circuit we report:
(1)~the consumed noise and whether it exceeds $\budget$ (a \emph{violation});
(2)~the execution time, measured with \todo{backend, e.g., SEAL BFV} on
\todo{hardware}, as the \todo{median} of \todo{$n$} runs;
and (3)~the number of distinct rotation keys and their serialized size in MB.
A serialized Galois key occupies approximately 18~MB in our setting
\todo{confirm for the parameters used here}.

% ---------------------------------------------------------------------
% Figure: per-benchmark comparison (style follows the CHEHAB RL paper:
% one bar group per benchmark). ILLUSTRATIVE DATA.
% ---------------------------------------------------------------------
\begin{figure*}[t]
  \centering
  \pgfplotsset{
    exp1bar/.style={
      ybar, bar width=4pt, width=\textwidth, height=3.6cm,
      enlarge x limits=0.04,
      symbolic x coords={DP16,DP32,L2-16,L2-32,HD16,HD32,PR8,LR8,BB4,MM4,SGx4,Max4,Sort3,PT},
      xtick=data, xticklabel style={font=\scriptsize},
      yticklabel style={font=\scriptsize}, ylabel style={font=\small},
      ymajorgrids, grid style={gray!25},
      legend style={font=\scriptsize, draw=none, legend columns=-1,
                    at={(0.5,1.03)}, anchor=south},
    }
  }
  \begin{tikzpicture}
  \begin{groupplot}[group style={group size=1 by 3, vertical sep=0.9cm,
                                 xticklabels at=edge bottom}]
  % (a) Noise consumption relative to the budget
  \nextgroupplot[exp1bar, ylabel={Noise / $\budget$}, ymin=0.7, ymax=1.25,
                 legend entries={\baseline{}, \cmorl{}}]
    \addplot[fill=gray!55, draw=gray!80] coordinates {
      (DP16,1.048) (DP32,1.122) (L2-16,1.026) (L2-32,1.139) (HD16,1.074)
      (HD32,1.178) (PR8,0.952) (LR8,0.817) (BB4,0.891) (MM4,1.096)
      (SGx4,0.930) (Max4,1.013) (Sort3,1.039) (PT,1.157)};
    \addplot[fill=blue!60, draw=blue!80] coordinates {
      (DP16,0.974) (DP32,0.987) (L2-16,0.961) (L2-32,0.996) (HD16,0.983)
      (HD32,0.991) (PR8,0.922) (LR8,0.817) (BB4,0.883) (MM4,0.983)
      (SGx4,0.913) (Max4,0.978) (Sort3,0.996) (PT,0.991)};
    \draw[red, dashed, thick] (rel axis cs:0,0.5454) -- (rel axis cs:1,0.5454)
      node[pos=0.0, above right, font=\scriptsize, text=red] {budget $\budget$};
  % (b) Speedup of CMO-RL over the baseline (execution time)
  \nextgroupplot[exp1bar, ylabel={Speedup over\\\baseline{}},
                 ylabel style={align=center}, ymin=0.9, ymax=1.1, bar width=6pt]
    \addplot[fill=blue!60, draw=blue!80] coordinates {
      (DP16,1.017) (DP32,1.022) (L2-16,0.998) (L2-32,1.023) (HD16,1.030)
      (HD32,1.063) (PR8,1.019) (LR8,1.000) (BB4,1.043) (MM4,1.028)
      (SGx4,1.017) (Max4,1.019) (Sort3,1.010) (PT,1.013)};
    \draw[black, thick] (rel axis cs:0,0.5) -- (rel axis cs:1,0.5);
  % (c) Rotation key size
  \nextgroupplot[exp1bar, ylabel={Rotation key\\size (MB)},
                 ylabel style={align=center}, ymin=0, ymax=140]
    \addplot[fill=gray!55, draw=gray!80] coordinates {
      (DP16,90) (DP32,108) (L2-16,90) (L2-32,108) (HD16,90) (HD32,126)
      (PR8,36) (LR8,36) (BB4,54) (MM4,72) (SGx4,54) (Max4,90) (Sort3,90) (PT,108)};
    \addplot[fill=blue!60, draw=blue!80] coordinates {
      (DP16,54) (DP32,72) (L2-16,54) (L2-32,72) (HD16,72) (HD32,72)
      (PR8,36) (LR8,18) (BB4,36) (MM4,54) (SGx4,36) (Max4,72) (Sort3,72) (PT,72)};
  \end{groupplot}
  \end{tikzpicture}
  \caption{\cmorl{} versus \baseline{} at noise budget
    $\budget=\ex{230}$~bits and preference $\vec{w}=(0.9,0.1)$.
    (a)~Consumed noise normalized by $\budget$; bars above the dashed line
    violate the budget.
    (b)~Execution-time speedup of \cmorl{} over \baseline{} (higher is
    better; values above~1 mean \cmorl{} is faster).
    (c)~Serialized rotation key size (lower is better).
    \ex{Illustrative placeholder data.}}
  \label{fig:exp1}
\end{figure*}

% ---------------------------------------------------------------------
% Table: per-benchmark details. ILLUSTRATIVE DATA.
% ---------------------------------------------------------------------
\begin{table*}[t]
  \centering
  \caption{Per-benchmark results at $\budget=\ex{230}$~bits and
    $\vec{w}=(0.9,0.1)$. Noise is the consumed noise in bits; \ding{55}
    marks a budget violation and \checkmark{} marks compliance. Time is the
    execution time in ms. Keys is the number of distinct rotation keys; MB
    is their serialized size. Best value per row in bold.
    \ex{Illustrative placeholder data.}}
  \label{tab:exp1-details}
  \small
  \setlength{\tabcolsep}{4pt}
  \begin{tabular}{@{}l rcrrr @{\hspace{12pt}} rcrrr@{}}
    \toprule
    & \multicolumn{5}{c}{\baseline{}} & \multicolumn{5}{c}{\cmorl{}} \\
    \cmidrule(lr){2-6}\cmidrule(l){7-11}
    Benchmark & Noise & & Time & Keys & MB & Noise & & Time & Keys & MB \\
    \midrule
    \ex{Dot Product ($N{=}16$)}       & \ex{241} & \ding{55} & \ex{412}  & \ex{5} & \ex{90}  & \ex{224} & \checkmark & \textbf{\ex{405}}  & \textbf{\ex{3}} & \textbf{\ex{54}} \\
    \ex{Dot Product ($N{=}32$)}       & \ex{258} & \ding{55} & \ex{655}  & \ex{6} & \ex{108} & \ex{227} & \checkmark & \textbf{\ex{641}}  & \textbf{\ex{4}} & \textbf{\ex{72}} \\
    \ex{L2 Distance ($N{=}16$)}       & \ex{236} & \ding{55} & \textbf{\ex{430}} & \ex{5} & \ex{90}  & \ex{221} & \checkmark & \ex{431}  & \textbf{\ex{3}} & \textbf{\ex{54}} \\
    \ex{L2 Distance ($N{=}32$)}       & \ex{262} & \ding{55} & \ex{790}  & \ex{6} & \ex{108} & \ex{229} & \checkmark & \textbf{\ex{772}}  & \textbf{\ex{4}} & \textbf{\ex{72}} \\
    \ex{Hamming Dist.\ ($N{=}16$)}    & \ex{247} & \ding{55} & \ex{980}  & \ex{5} & \ex{90}  & \ex{226} & \checkmark & \textbf{\ex{951}}  & \textbf{\ex{4}} & \textbf{\ex{72}} \\
    \ex{Hamming Dist.\ ($N{=}32$)}    & \ex{271} & \ding{55} & \ex{1490} & \ex{7} & \ex{126} & \ex{228} & \checkmark & \textbf{\ex{1402}} & \textbf{\ex{4}} & \textbf{\ex{72}} \\
    \ex{Poly.\ Regression ($N{=}8$)}  & \ex{219} & \checkmark & \ex{315} & \ex{2} & \ex{36}  & \ex{212} & \checkmark & \textbf{\ex{309}}  & \ex{2} & \ex{36} \\
    \ex{Lin.\ Regression ($N{=}8$)}   & \ex{188} & \checkmark & \ex{140} & \ex{2} & \ex{36}  & \ex{188} & \checkmark & \ex{140}  & \textbf{\ex{1}} & \textbf{\ex{18}} \\
    \ex{Box Blur ($N{=}4$)}           & \ex{205} & \checkmark & \ex{96}  & \ex{3} & \ex{54}  & \ex{203} & \checkmark & \textbf{\ex{92}}   & \textbf{\ex{2}} & \textbf{\ex{36}} \\
    \ex{Matrix Mul.\ ($N{=}4$)}       & \ex{252} & \ding{55} & \ex{512}  & \ex{4} & \ex{72}  & \ex{226} & \checkmark & \textbf{\ex{498}}  & \textbf{\ex{3}} & \textbf{\ex{54}} \\
    \ex{Sobel Gx ($N{=}4$)}           & \ex{214} & \checkmark & \ex{530} & \ex{3} & \ex{54}  & \ex{210} & \checkmark & \textbf{\ex{521}}  & \textbf{\ex{2}} & \textbf{\ex{36}} \\
    \ex{Maximum ($N{=}4$)}            & \ex{233} & \ding{55} & \ex{1210} & \ex{5} & \ex{90}  & \ex{225} & \checkmark & \textbf{\ex{1188}} & \textbf{\ex{4}} & \textbf{\ex{72}} \\
    \ex{Sort ($N{=}3$)}               & \ex{239} & \ding{55} & \ex{1380} & \ex{5} & \ex{90}  & \ex{229} & \checkmark & \textbf{\ex{1366}} & \textbf{\ex{4}} & \textbf{\ex{72}} \\
    \ex{Poly.\ Tree (100/50, $d{=}5$)}& \ex{266} & \ding{55} & \ex{2010} & \ex{6} & \ex{108} & \ex{228} & \checkmark & \textbf{\ex{1985}} & \textbf{\ex{4}} & \textbf{\ex{72}} \\
    \bottomrule
  \end{tabular}
\end{table*}

% ---------------------------------------------------------------------
% Table: aggregate summary. ILLUSTRATIVE DATA.
% ---------------------------------------------------------------------
\begin{table}[t]
  \centering
  \caption{Aggregate results over the \ex{14} benchmarks of
    \cref{tab:exp1-details}. Geometric means are taken over per-benchmark
    ratios relative to \baseline{}. \ex{Illustrative placeholder data.}}
  \label{tab:exp1-summary}
  \small
  \begin{tabular}{@{}lrr@{}}
    \toprule
    Metric & \shortstack{\chehabrl{}\\+\rks{}} & \cmorl{} \\
    \midrule
    Budget violations                         & \ex{10/14} & \textbf{\ex{0/14}} \\
    Max.\ overshoot (\% of $\budget$)& \ex{17.8\%} & \textbf{\ex{0\%}} \\
    Geomean speedup& $1.00\times$ & \textbf{\ex{1.02$\times$}} \\
    Lower-or-equal time& \ex{2/14} & \textbf{\ex{13/14}} \\
    Total rotation keys                       & \ex{64} & \textbf{\ex{44}} \\
    Total key size (MB)& \ex{1152} & \textbf{\ex{792}} \\
    Geomean key-size ratio& $1.00$ & \textbf{\ex{0.69}} \\
    \bottomrule
  \end{tabular}
\end{table}

\paragraph{Constraint satisfaction.}
\cref{fig:exp1}(a) and \cref{tab:exp1-details} show that \baseline{} exceeds
the noise budget on \ex{10} of the \ex{14} benchmarks, whereas \cmorl{}
satisfies the budget on all \ex{14}. The largest violation of \baseline{}
occurs on \ex{Hamming Distance ($N{=}32$)}, whose circuit consumes \ex{271}
bits, \ex{17.8\%} above the budget. A violation is not a small loss of
quality. Under fixed encryption parameters, a circuit that consumes more
noise than the budget allows cannot be guaranteed to decrypt correctly; the
only remedy is to choose larger encryption parameters, which increases both
execution time and key size for the whole program.
\todo{Confirm with measured SEAL remaining-noise budgets or decryption
checks that the violating circuits indeed fail, or that larger parameters
are required.}
The four benchmarks on which \baseline{} complies (\ex{Polynomial
Regression, Linear Regression, Box Blur, and Sobel Gx}) have shallow
circuits that consume at most \ex{219} bits under either system. For these
benchmarks the budget is not binding, and \cmorl{} behaves as an
unconstrained co-optimizer of time and memory.

\paragraph{Execution time.}
\cref{fig:exp1}(b) shows that satisfying the budget does not cost \cmorl{}
execution time. \cmorl{} is faster than or as fast as \baseline{} on
\ex{13} of the \ex{14} benchmarks, with a geometric-mean speedup of
\ex{$1.02\times$} (range \ex{$0.998\times$}--\ex{$1.063\times$};
\cref{tab:exp1-summary}). The one benchmark on which \cmorl{} is slower,
\ex{L2 Distance ($N{=}16$)}, differs by \ex{1~ms (0.2\%)}, which is
\todo{within the run-to-run variation of the measurements; report the
standard deviation}. The largest speedup, \ex{$1.063\times$}, appears on
\ex{Hamming Distance ($N{=}32$)}, the benchmark where \rks{} decomposes the
most rotations.
We attribute this speedup to the order in which the two systems handle key
memory. \baseline{} first produces a speed-optimized circuit and then
replaces infrequent rotations with chains of power-of-two rotations; each
chain executes several rotations where the original circuit executed one.
\cmorl{} reduces key memory while it chooses the vectorization, so it can
select data layouts that reuse a small set of rotation offsets and need no
rotation chains.
\todo{Support this explanation with operation counts: number of rotations
inserted by \rks{} in \baseline{} vs.\ total rotations in \cmorl{}.}

\paragraph{Rotation key memory.}
\cref{fig:exp1}(c) shows that \cmorl{} requires fewer or equally many
rotation keys on every benchmark. Over all \ex{14} benchmarks, \cmorl{}
uses \ex{44} keys (\ex{792}~MB) against \ex{64} keys (\ex{1152}~MB) for
\baseline{}, a reduction of \ex{31\%} in total key size; the
geometric-mean key-size ratio is \ex{0.69} (\cref{tab:exp1-summary}). The
reduction is obtained even though \baseline{} applies a dedicated
key-reduction pass and \cmorl{} assigns only a weight of $\wkeys = 0.1$ to
memory. On \ex{Polynomial Regression ($N{=}8$)}, both systems require the
same \ex{2} keys, which suggests that \ex{2} offsets are necessary for this
circuit under both optimization strategies.

\paragraph{Why \baseline{} violates the budget.}
The noise model shows which operations consume the noise in the violating
circuits. \todo{Replace this paragraph with the measured breakdown; the
explanation below is a hypothesis to be confirmed.}
\chehabrl{} vectorizes aggressively because vector operations reduce its
execution-cost model. Vectorization introduces rotations and plaintext
masking multiplications, and in BFV each of these operations adds noise.
\rks{} then lengthens the rotation sequences on the critical path by
replacing single rotations with chains of rotations, which adds further
key-switching noise. Neither step observes the noise budget. \cmorl{}
receives the noise budget and the current noise ratio as part of its state
and is penalized when the circuit exceeds the budget, so it learns to stop
applying noise-heavy rewrites once the remaining margin is small.

\paragraph{Summary.}
At a noise budget of \ex{230}~bits and a preference of $\vec{w}=(0.9,0.1)$,
\cmorl{} satisfies the budget on all \ex{14} benchmarks, whereas \baseline{}
violates it on \ex{10}. \cmorl{} achieves this with equal or lower execution
time on \ex{13} benchmarks (geometric-mean speedup \ex{$1.02\times$}) and
with \ex{31\%} less rotation key memory in total. These results hold for one
budget and one preference; \cref{sec:exp2} \todo{check reference} evaluates
how \cmorl{} behaves across budgets and preferences.
